Low-count positron emission tomography (PET) is degraded by signal-dependent Poisson fluctuations, whereas indiscriminate smoothing can erase small lesions. This study proposes Adaptive VST-PGD, an image-preprocessing method that combines the Anscombe variance-stabilizing transform with non-negative, spatially adaptive first- and second-order regularization. The resulting objective is minimized by projected gradient descent using a Barzilai--Borwein spectral step and Armijo backtracking. Experiments on a reproducible $256\times256$ PET-like Shepp--Logan phantom compare the method with Gaussian filtering, non-local means, and total-variation gradient descent at 50\%, 25\%, and 10\% dose levels. Adaptive VST-PGD achieves the highest PSNR at all three levels (\ProposedPSNRFifty, \ProposedPSNRTwentyFive, and \ProposedPSNRTen{} dB, respectively). Ablation, convergence, and parameter-sensitivity analyses quantify the contributions of variance stabilization, adaptive weighting, and spectral step selection. These results establish an interpretable numerical baseline for low-dose PET preprocessing; clinical effectiveness remains to be validated on physical phantom and patient data.
